% Koordinaten für den Hintergrundrahmen
\newcommand{\diameter}{20}
\newcommand{\xone}{-15}
\newcommand{\xtwo}{160}
\newcommand{\yone}{15}
\newcommand{\ytwo}{-253}

\newcommand{\hoehea}{60}
\newcommand{\hoeheb}{60}




\begin{titlepage}
	% Hintergrundrahmen
	\begin{tikzpicture}[overlay]
	\draw[color=gray]  
			(\xone mm, \yone mm)
	  -- (\xtwo mm, \yone mm)
	arc (90:0:\diameter pt) 
	  -- (\xtwo mm + \diameter pt , \ytwo mm) 
		-- (\xone mm + \diameter pt , \ytwo mm)
	arc (270:180:\diameter pt)
		-- (\xone mm, \yone mm);
	\end{tikzpicture}
	
	% KIT-Logo
	\begin{textblock}{10}[0,0](4.5,2.5)
		\includegraphics[width=.25\textwidth]{include/kitlogo.pdf}
	\end{textblock}
	\changefont{phv}{m}{n}	% helvetica
	\begin{textblock}{10}[0,0](5.5,2.2)
		\begin{flushright}
			\Large FAKULTÄT FÜR PHYSIK\\Praktikum Klassische Physik
		\end{flushright}
	\end{textblock}
	
	\begin{textblock}{10}[0,0](4.2,3.1)
		\begin{tikzpicture}[overlay]
		\draw[color=gray]
				(\xone mm + 5 mm, -12 mm)
		  -- (\xtwo mm + \diameter pt - 5 mm, -12 mm);
		\end{tikzpicture}
	\end{textblock}
	
	\Large
	% Zeile 1
	\begin{textblock}{12}[0,0](3.85,3.85)
		\mytextfield{Prak.}{\praktikum}{0.9cm}{17pt}{P1/P2}{2}{Praktikum}
	\end{textblock}
	\begin{textblock}{12}[0,0](5.8,3.85)
		\mytextfield{Semester}{\semester}{2.6cm}{17pt}{z.B. \glqq WS14/15\grqq\ oder \glqq SS15\grqq}{0}{Semester}
	\end{textblock}
	\begin{textblock}{12}[0,0](9.8,3.85)
		\mytextfield{Wochentag}{\wochentag}{1.3cm}{17pt}{Mo/Di/Mi/Do}{2}{Wochentag}
	\end{textblock}
	\begin{textblock}{12}[0,0](13.15,3.85)
		\mytextfield{Gruppennr.}{\gruppennr}{1.06cm}{17pt}{\#\#}{2}{Gruppennummer}
	\end{textblock}
	
	% Zeile 2
	\begin{textblock}{12}[0,0](3.85,4.4)
		\mytextfield{Name}{\nachnamea}{6cm}{17pt}{}{0}{Name1}
	\end{textblock}
	\begin{textblock}{12}[0,0](9.8,4.4)
		\mytextfield{Vorname}{\vornamea}{6cm}{17pt}{}{0}{Vorname1}
	\end{textblock}
	
	% Zeile 3
	\begin{textblock}{12}[0,0](3.85,4.95)
		\mytextfield{Name}{\nachnameb}{6cm}{17pt}{}{0}{Name2}
	\end{textblock}
	\begin{textblock}{12}[0,0](9.8,4.95)
		\mytextfield{Vorname}{\vornameb}{6cm}{17pt}{}{0}{Vorname2}
	\end{textblock}
	
	% Zeile 4
	\begin{textblock}{12}[0,0](3.85,5.5)
		\normalsize\mytextfield{Emailadressen}{fachschaft@physik.kit.edu}{13.41cm}{10pt}{Optional}{0}{Emailadressen}
	\end{textblock}
	
	% Zeile 5
	\begin{textblock}{12}[0,0](3.85,6.45)
		\mytextfield{Versuch}{\versuch\ (\praktikum-\versuchsnr)}{9.45cm}{14pt}{z.B. \glqq Galvanometer (P1-13)\grqq\ oder \glqq Mikrowellenoptik (P2-15)\grqq}{0}{Versuch}
	\end{textblock}
	\begin{textblock}{12}[0,0](12.85,6.45)
		\mytextfield{Fehlerrech.}{\fehlerrechnung}{1.46cm}{17pt}{Ja/Nein}{4}{Fehlerrechnung}
	\end{textblock}
	
	% Zeile 6
	\begin{textblock}{12}[0,0](3.85,7)
		\mytextfield{Betreuer}{\betreuer}{7cm}{17pt}{}{0}{Betreuer}
	\end{textblock}
	\begin{textblock}{12}[0,0](11.09,7)
		\mytextfield{Durchgeführt am}{\durchgefuehrt}{2.53cm}{17pt}{TT.MM.JJ}{8}{Durchfuehrung}
	\end{textblock}
	
	%querstrich
	\begin{textblock}{20}[0,0](0,7.9)\tiny\centering
		Wird vom Betreuer ausgefüllt.
	\end{textblock}
	\begin{tikzpicture}[overlay]
	\draw[color=gray]
			(\xone mm + 5 mm, -95 mm)
	  -- (\xtwo mm + \diameter pt - 5 mm, -95 mm);
	\end{tikzpicture}
	
	% Zeile 1
	\begin{textblock}{12}[0,0](3.85,8.27)
		\myTtextfield{1. Abgabe am}{}{2.5cm}{17pt}{}
	\end{textblock}
	
	% Block 1
	\begin{tikzpicture}[overlay]
	\draw[color=gray]  
			(\xone mm + 10 mm, -107.5 mm)
	  -- (\xtwo mm + \diameter pt - 10 mm, -107.5 mm)
	  -- (\xtwo mm + \diameter pt - 10 mm, -107.5 mm - \hoehea mm)
	  -- (\xone mm + 10 mm, -107.5 mm - \hoehea mm)
	  -- (\xone mm + 10 mm, -107.5 mm);
	\end{tikzpicture}
	\begin{textblock}{12}[0,0](4,9)
		\myTtextfield{Rückgabe am}{}{2.5cm}{17pt}{}
	\end{textblock}
	\begin{textblock}{12}[0,0](8.6,9.1)
		Begründung:
	\end{textblock}
	
	% Zeile 2
	\begin{textblock}{12}[0,0](3.85,12.3)
		\myTtextfield{2. Abgabe am}{}{2.5cm}{17pt}{}
	\end{textblock}
	
	% Block 2
	\begin{tikzpicture}[overlay]
	\draw[color=gray]  
			(\xone mm + 10 mm, -180 mm)
	  -- (\xtwo mm + \diameter pt - 10 mm, -180 mm)
	  -- (\xtwo mm + \diameter pt - 10 mm, -180 mm - \hoehea mm)
	  -- (\xone mm + 10 mm, -180 mm - \hoehea mm)
	  -- (\xone mm + 10 mm, -180 mm);
	\end{tikzpicture}
	\begin{textblock}{12}[0,0](4,13.2)
		Ergebnis:~~~~+~~~/~~~0~~~/~~~-
	\end{textblock}
	\begin{textblock}{12}[0,0](9.5,13.2)
		Fehlerrechnung:~~~Ja~~~/~~~Nein
	\end{textblock}
	\begin{textblock}{12}[0,0](4,13.6)
		\myTtextfield{Datum}{}{2.5cm}{17pt}{}
	\end{textblock}
	\begin{textblock}{12}[0,0](8.5,13.6)
		\myTtextfield{Handzeichen}{}{5.5cm}{17pt}{}
	\end{textblock}
	\begin{textblock}{12}[0,0](4,14.2)\Large
		Bemerkungen:
	\end{textblock}
	
	
	
	% unterste Zeile über das KIT
	\begin{textblock}{10}[0,0](4,16.8)
		\tiny{KIT -- Universität des Landes Baden-Württemberg und nationales Forschungszentrum in der Helmholtz-Gemeinschaft}
	\end{textblock}
	\begin{textblock}{10}[0,0](14,16.75)
		\large{\textbf{www.kit.edu}}
	\end{textblock}
\end{titlepage}
